\documentclass{article}

\usepackage[preprint,nonatbib]{neurips_2023}
\usepackage[T1]{fontenc}
\usepackage[utf8]{inputenc}
\usepackage{microtype}
\usepackage{booktabs}
\usepackage{amsmath}
\usepackage[colorlinks=true,allcolors=blue]{hyperref}

\title{Content-Dependent Acoustic Memory}
\author{CleaRx}

\begin{document}
\maketitle

\section{Question}

\textbf{Under a fixed memory budget, can Qwen3-Omni preserve more useful speech history by allocating acoustic positions according to the information they add beyond the transcript?} \href{https://arxiv.org/html/2608.08569v1}{VoxZip}, the immediate prior method, pools each ASR-defined interval to its transcript-token count, fuses audio and text by addition, then prunes KV entries~\hyperlink{ref:voxzip}{[1]}. Its acoustic budget is therefore fixed by transcript length.

We instead allocate \(b_i=g_\phi(z_i,B)\) positions to segment \(i\) according to evidence \(z_i\), subject to \(\sum_i b_i\leq B\). The budget should follow \emph{marginal acoustic value}: evidence not already carried by text.

\section{Preliminary evidence}

Three frozen-Qwen3-Omni experiments support this question~\hyperlink{ref:clearx}{[3]}. Mean-pooled audio reduced teacher KL from 2.136 to 1.966 in one clip and improved mean KL by 7\% across ten clips, with gains on 7/10. Continuous weighting was 5.8\% worse than the best fixed weight and saved no positions. Under a fixed 48-position budget, agreement-based retention improved KL by 7.96\% over uniform retention, but not an audio-only allocator. Its robust contribution was the measured heterogeneity:

\begin{center}
\small
\begin{tabular}{lrr}
\toprule
Benefit of four positions over transcript-only & Clean & Degraded \\
\midrule
Exact transcript & 0.0025 & 0.0011 \\
Incomplete transcript & 4.5444 & 1.7199 \\
Conflicting transcript & 4.3288 & 0.8755 \\
\bottomrule
\end{tabular}
\end{center}

Exact transcripts gained almost nothing; incomplete or conflicting transcripts gained much more, especially with clean audio.

\section{Study}

\textbf{1. Allocation and representation.} In multi-turn histories, measure utility at \(b_i\in\{0,1,2,4,8\}\) and allocate one shared budget. At equal KV bytes, compare uniform, random, audio-only, nonlinear, and oracle allocation; separately compare mean pooling, VoxZip addition, learned summarization, and learned fusion. This isolates where memory goes from what retained positions contain.

\textbf{2. Compression location.} Compare the best pre-Thinker method with SpeechKV-style pooling inside early and middle Thinker layers~\hyperlink{ref:speechkv}{[2]}. With Qwen3-Omni frozen, measure KL, task agreement, KV memory, and latency on held-out conversations.

Adaptive allocation must beat uniform, random, and audio-only policies at equal cost; learned fusion must beat VoxZip-style addition. The result will map what acoustic evidence merits retention, how much memory it receives, and where it can first be compressed safely.

{\small\paragraph{References.}
\hypertarget{ref:voxzip}{[1]} Liu et al. \href{https://arxiv.org/html/2608.08569v1}{VoxZip}. ACM MM, 2026.
\quad \hypertarget{ref:speechkv}{[2]} Kim et al. \href{https://arxiv.org/html/2607.06827v1}{SpeechKV}. 2026.
\quad \hypertarget{ref:clearx}{[3]} CleaRx. \href{https://github.com/pvd232/CleaRx/tree/main/experiments}{Exploratory audio-memory results}. 2026.}

\end{document}
